\chapter{Background work}

Background jobs move slow or unreliable work out of the request path.
Use them for emails, reports, and heavy processing.

\section{Queues and workers}
A queue decouples producers from consumers.
Workers pull jobs and retry failures.

\section{Retries}
Retries should be bounded and idempotent.
Always include a backoff strategy to avoid overload.

\section{Dead letter queues}
A dead letter queue captures jobs that consistently fail.
It is the safest place to debug without blocking the main system.

\section{Job idempotency}
Use a job key so the same job can run twice without harm.
This is critical during deploys and retries.

\section{When to use background work}
Background jobs are best for work that is slow, unreliable, or can be delayed.
Do not move user-critical work to a queue just to avoid fixing latency.
If the user expects immediate results, keep a synchronous path.

\section{Queue types}
Queues come in multiple forms:
\begin{itemize}[nosep]
  \item FIFO queues for ordered processing
  \item standard queues for throughput
  \item delayed queues for scheduled retries
\end{itemize}
Pick the queue type that matches ordering and latency requirements.

\section{At-least-once delivery}
Most queues deliver at least once.
Your job handlers must tolerate duplicates.
Idempotency is non-negotiable in this model.

\section{Visibility timeouts}
Visibility timeouts define how long a job is locked by a worker.
If the timeout is too short, jobs run twice.
If it is too long, failed jobs block progress.

\section{Retry policy}
Retries should be explicit and bounded.
Use exponential backoff with jitter.
A typical retry policy uses 3 to 5 attempts with increasing delays.

\section{Retry schedule example}
\begin{lstlisting}[language=YAML]
retry_policy:
  max_attempts: 5
  backoff_seconds: [5, 30, 120, 300, 900]
  jitter: true
\end{lstlisting}

\section{Dead letter handling}
Dead letter queues should be monitored.
Define a process for triage, not just storage.
Automate replays when safe, and require manual review for risky jobs.

\section{Poison pill jobs}
Some jobs fail because of data corruption or invalid state.
These should be marked as poison pills and quarantined.
Do not retry forever; it wastes capacity.

\section{Job payload design}
Keep payloads small and explicit.
Include identifiers, not entire objects.
Large payloads increase queue cost and reduce flexibility.

\section{Priority queues}
Not all jobs are equal.
Use priorities to keep user-facing work ahead of batch processing.
If you mix priorities, document the policy clearly.

\section{Scheduling and cron}
Scheduled work should be expressed as recurring jobs.
Use a dedicated scheduler rather than ad-hoc timers.
Schedulers should support idempotent retries like any other job.

\section{Rate limiting jobs}
External APIs often impose rate limits.
Use token buckets or per-provider queues to avoid throttling.
Rate limiting inside the worker prevents accidental abuse.

\section{Concurrency control}
Workers should bound concurrency per job type.
A single job class can starve others if not controlled.
Use separate queues or worker pools when necessary.

\section{Ordering requirements}
If order matters, use FIFO queues and deduplication keys.
Ordering reduces throughput, so only require it when needed.
Document which jobs require ordering.

\section{Idempotency keys}
Store idempotency keys with a short TTL.
Keys must cover all side effects of the job.
If a retry happens, reusing the key should result in a no-op.

\section{Example job record}
\begin{lstlisting}[language=JSON]
{
  "job_id": "job_123",
  "type": "send_email",
  "payload": {"user_id": "user_7"},
  "idempotency_key": "email_user_7_welcome",
  "attempt": 2
}
\end{lstlisting}

\section{Backpressure}
If the queue grows too fast, slow down producers.
Backpressure keeps the system stable during spikes.
Expose queue depth as a metric to producers.

\section{Scaling workers}
Scale workers based on queue depth and processing time.
Fast jobs need more workers than slow jobs.
Autoscaling should consider both backlog and throughput.

\section{Job timeouts}
Set execution timeouts for jobs.
If a job exceeds its timeout, fail it and retry if safe.
Timeouts prevent long-running jobs from blocking workers.

\section{Cancellation}
Some jobs should be cancelable.
Support cancellation by checking a cancel flag at safe points.
Cancellation is valuable for user-driven actions.

\section{Bulk processing}
Bulk jobs should process in chunks.
Chunking allows partial progress and easier retries.
Track progress and resume from the last successful chunk.

\section{Observability for jobs}
Record job metrics:
\begin{itemize}[nosep]
  \item queue depth
  \item job latency
  \item failure rate
  \item retries per job
\end{itemize}
These metrics reveal saturation and stability issues.

\section{Example job metrics}
\begin{lstlisting}[language=YAML]
metrics:
  queue_depth: 1200
  job_latency_p95_ms: 800
  retry_rate: 0.07
\end{lstlisting}

\section{Security considerations}
Jobs often handle sensitive data.
Do not log full payloads.
Use scoped credentials for workers and rotate them regularly.

\section{Consistency with the request path}
Background jobs should not bypass business rules.
Reuse core domain logic where possible.
If background logic diverges, data corruption is likely.

\section{Job lifecycle}
Jobs should have clear states: queued, running, succeeded, failed.
Persist state transitions for debugging and audits.
State visibility makes operations predictable.

\section{Leasing and locks}
Workers should acquire a lease when they start a job.
Leases prevent multiple workers from running the same job.
If a worker dies, the lease expires and the job can be retried.

\section{Exactly-once illusion}
Exactly-once processing is rare in practice.
Aim for at-least-once with idempotency.
If you need exactly-once outcomes, use deduplication and idempotent writes.

\section{Deduplication storage}
Deduplication keys need storage.
Use a small table or cache keyed by idempotency key.
Set TTLs based on retry windows and business rules.

\section{Payload versioning}
Job payloads evolve.
Include a \code{payload\_version} field and handle old versions.
Versioning prevents deploys from breaking queued jobs.

\section{Schema validation}
Validate job payloads before enqueueing.
Reject invalid jobs early to avoid poisoned queues.
Schema validation should be shared between producers and consumers.

\section{Job state table}
\begin{tabularx}{\textwidth}{@{}lX@{}}
\toprule
State & Meaning \\
\midrule
queued & waiting to be processed \\
running & currently being handled \\
failed & exceeded retry policy \\
succeeded & completed without error \\
\bottomrule
\end{tabularx}

\section{Fanout and chaining}
Fanout jobs create multiple downstream jobs.
Chaining jobs ensures a sequence of dependent steps.
Both patterns need clear error handling and rollbacks.

\section{Workflow orchestration}
Complex workflows benefit from orchestration.
Use a workflow engine if you need stateful retries and branching.
Ad-hoc workflows grow fragile quickly.

\section{Backfills}
Backfills are large background jobs that rebuild data.
Run backfills in controlled batches to avoid overload.
Track progress and support pause and resume.

\section{Reprocessing policy}
Sometimes you need to replay jobs after bugs.
Store enough metadata to reprocess safely.
Document which job types can be replayed.

\section{Priority inversion}
Priority queues can suffer inversion.
If low-priority work blocks high-priority jobs, split queues.
Separate worker pools reduce interference.

\section{Job sharding}
Shard heavy jobs by tenant or region.
Sharding improves parallelism and reduces contention.
Shards should align with data locality when possible.

\section{Storage considerations}
Job queues are not long-term storage.
If you need long retention, persist job metadata elsewhere.
Queues should be treated as transient buffers.

\section{Operational tooling}
Operators need tools to inspect and requeue jobs.
Provide a safe admin interface for job operations.
Manual fixes should be logged and auditable.

\section{Example job admin actions}
\begin{lstlisting}[language=YAML]
admin_actions:
  - replay_job
  - cancel_job
  - move_to_dlq
\end{lstlisting}

\section{Job SLAs}
Background work still has latency expectations.
Define an SLA for job completion.
If jobs exceed the SLA, it should trigger investigation.

\section{Job latency tracking}
Track end-to-end job latency from enqueue to completion.
Separate queue wait time from processing time.
This helps determine whether the bottleneck is capacity or code.

\section{Multi-tenant fairness}
Large tenants can monopolize queues.
Use per-tenant rate limits or weighted queues.
Fairness prevents smaller tenants from being starved.

\section{Routing by job type}
Separate queues by job type when workloads differ.
CPU-bound and IO-bound jobs should not share the same pool.
Routing simplifies tuning and scaling.

\section{Blue/green for workers}
Worker deploys can cause job failures.
Use blue/green or rolling deploys for worker fleets.
Drain jobs before terminating instances when possible.

\section{Graceful shutdown}
Workers should finish the current job before shutting down.
Respect the visibility timeout and release leases cleanly.
Abrupt termination causes retries and duplicates.

\section{Poison pill workflow}
Define a workflow for poison pills:
\begin{itemize}[nosep]
  \item quarantine the job
  \item inspect payload and logs
  \item patch data or logic
  \item replay safely
\end{itemize}

\section{Example job router}
\begin{lstlisting}[language=Python]
def route(job):
    if job.type in {"send_email", "send_sms"}:
        return "notifications"
    if job.type in {"generate_report"}:
        return "reports"
    return "default"
\end{lstlisting}

\section{Backfill safety}
Backfills should be resumable.
Store checkpoints after each batch.
If a backfill fails, resume rather than restarting.

\section{Costs and limits}
Queues and workers have cost profiles.
Track cost per job type.
Cost visibility helps prioritize optimizations.

\section{Side effects}
Jobs often trigger side effects like emails or external API calls.
Design side effects to be idempotent or safe to repeat.
If the side effect is not idempotent, store a side-effect record first.

\section{Checkpointing}
Long-running jobs should checkpoint progress.
Checkpointing allows retries to resume rather than restart.
Store checkpoints in durable storage.

\section{Batch windowing}
Batch processing benefits from windowing.
A small window reduces latency, a large window reduces overhead.
Choose a window size based on user expectations.

\section{Dedupe windows}
Deduplication should have a window.
A short window reduces storage cost but allows rare duplicates.
A long window increases cost but improves safety.

\section{Example dedupe record}
\begin{lstlisting}[language=JSON]
{
  "idempotency_key": "email_user_7_welcome",
  "created_at": "2026-02-18T10:00:00Z",
  "expires_at": "2026-02-25T10:00:00Z"
}
\end{lstlisting}

\section{Alerts for queues}
Alert on queue depth and job age.
Queue depth indicates capacity issues.
Job age indicates latency and SLA risk.

\section{Example queue alert}
\begin{lstlisting}[language=YAML]
alert: QueueDepthHigh
expr: queue_depth > 5000
for: 10m
labels:
  severity: "sev2"
\end{lstlisting}

\section{Job SLAs table}
\begin{tabularx}{\textwidth}{@{}lXr@{}}
\toprule
Job type & Purpose & SLA (min) \\
\midrule
send\_email & user notifications & 5 \\
generate\_report & analytics export & 30 \\
reindex\_search & rebuild index & 120 \\
\bottomrule
\end{tabularx}

\section{Payload size}
Large payloads slow queues and increase cost.
Compress payloads only when necessary.
Prefer storing large objects in a database and referencing them by id.

\section{Job testing}
Background jobs should have unit and integration tests.
Test idempotency and retry behavior explicitly.
Jobs are part of your production surface and deserve the same rigor.

\section{Schema registry}
If you have many job types, use a schema registry.
Schema registry helps validate payloads across teams.
It prevents breaking changes to job contracts.

\section{Reconciliation jobs}
Reconciliation jobs fix drift between systems.
Run them on a schedule and log corrections.
Reconciliation reduces long-term data inconsistency.

\section{Outbox-driven jobs}
Use outbox tables to trigger background work reliably.
The outbox ensures jobs are not lost during failures.
Outbox-driven jobs reduce dual-write risks.

\section{Partitioning queues}
Partition queues by key for fairness.
Partitions allow parallelism without losing ordering.
Choose partition keys based on tenant or resource id.

\section{Example partition config}
\begin{lstlisting}[language=YAML]
queue:
  partitions: 8
  partition_key: "tenant_id"
\end{lstlisting}

\begin{checklistbox}{Checklist: background work}
\begin{itemize}[nosep]
  \item Keep jobs idempotent.
  \item Use retries with backoff.
  \item Monitor queue depth.
\end{itemize}
\end{checklistbox}
